\documentclass[11pt]{article}

\usepackage{amsmath, amssymb, amsthm}

\title{Constraint-Gated Decision Systems and Threshold Stability}
\author{}
\date{}

\theoremstyle{plain}
\newtheorem{theorem}{Theorem}
\newtheorem{proposition}{Proposition}

\theoremstyle{definition}
\newtheorem{definition}{Definition}

\begin{document}

\maketitle

\begin{abstract}
We introduce a constraint-gated decision framework in which actions are permitted only when a normalized alignment score exceeds a fixed threshold. The framework is defined in terms of inner products in a finite-dimensional space and admits a natural geometric interpretation. We show that thresholds bounded away from zero induce stable decision regions, while vanishing thresholds admit degenerate behavior. A canonical threshold at angle $\pi/4$ yields a normalized bound $1/\sqrt{2}$, providing a symmetric reference point for balanced decision criteria.
\end{abstract}

\section{Introduction}

Decision systems often rely on scalar scores to determine whether an action should be taken. We consider a simple model in which such decisions are governed by a constraint threshold applied to a normalized alignment score.

The purpose of this paper is to define this mechanism in a minimal mathematical setting and to analyze the stability properties induced by a fixed threshold.

\section{Model}

Let $V = \mathbb{R}^d$ be a finite-dimensional real vector space with standard inner product $\langle \cdot, \cdot \rangle$.

\begin{definition}
Let $x \in V$ be a signal vector and $c \in V$ be a constraint vector with $\|c\| = 1$. Define the alignment score
\[
s(x) = \frac{\langle x, c \rangle}{\|x\|},
\]
whenever $x \neq 0$.
\end{definition}

Thus $s(x)$ is the cosine of the angle between $x$ and $c$, and satisfies
\[
-1 \le s(x) \le 1.
\]

\begin{definition}
Fix a threshold $\tau \in [0,1]$. The constraint gate is the decision rule:
\[
G_\tau(x) =
\begin{cases}
\text{allow} & \text{if } s(x) \ge \tau, \\
\text{block} & \text{otherwise}.
\end{cases}
\]
\end{definition}

\section{Geometric Interpretation}

The condition $s(x) \ge \tau$ is equivalent to
\[
\cos \theta \ge \tau,
\]
where $\theta$ is the angle between $x$ and $c$.

Thus, the allowed region is a cone:
\[
\theta \le \arccos(\tau).
\]

\section{Stability of Thresholds}

We now characterize how the threshold $\tau$ affects the decision region.

\begin{proposition}
If $\tau > 0$, then the allowed region excludes all vectors orthogonal to $c$.
\end{proposition}

\begin{proof}
If $x$ is orthogonal to $c$, then $\langle x, c \rangle = 0$, so $s(x)=0 < \tau$.
\end{proof}

\begin{theorem}
Let $\tau > 0$. Then the allowed region
\[
\{ x \in V : s(x) \ge \tau \}
\]
is a closed convex cone with nonzero angular separation from its complement.
\end{theorem}

\begin{proof}
The set $\{x : \langle x, c \rangle \ge \tau \|x\|\}$ is closed and convex. The condition $\tau > 0$ implies a strictly positive angular gap between allowed and blocked regions.
\end{proof}

\section{Degeneracy at Zero Threshold}

\begin{theorem}
If $\tau = 0$, then the gate reduces to
\[
G_0(x) = \text{allow} \quad \text{if and only if} \quad \langle x, c \rangle \ge 0.
\]
In this case, the boundary between allowed and blocked regions has no angular separation.
\end{theorem}

\begin{proof}
If $\tau = 0$, then $s(x) \ge 0$ is equivalent to $\langle x, c \rangle \ge 0$. The separating hyperplane is $\langle x, c \rangle = 0$, which has zero angular margin.
\end{proof}

Thus, $\tau = 0$ yields a degenerate decision boundary.

\section{Reference Threshold at $\pi/4$}

A symmetric threshold is obtained by setting
\[
\tau = \frac{1}{\sqrt{2}},
\]
which corresponds to angle
\[
\theta = \frac{\pi}{4}.
\]

\begin{proposition}
If $\tau = 1/\sqrt{2}$, then the allowed region consists of all vectors forming an angle at most $\pi/4$ with $c$.
\end{proposition}

\begin{proof}
Immediate from $\cos(\pi/4)=1/\sqrt{2}$.
\end{proof}

This threshold provides a balanced criterion: the projection of $x$ onto $c$ must be at least a fixed fraction of its norm.

\section{Robustness Under Perturbation}

\begin{proposition}
Let $\tau > 0$. Then there exists $\varepsilon > 0$ such that if $x$ satisfies $s(x) \ge \tau + \varepsilon$, small perturbations of $x$ do not change the decision outcome.
\end{proposition}

\begin{proof}
Continuity of $s(x)$ implies stability under perturbations away from the threshold boundary.
\end{proof}

Thus, positive thresholds create a buffer region that stabilizes decisions.

\section{Conclusion}

We have defined a constraint-gated decision system based on normalized alignment scores. The key structural distinction is:

\begin{itemize}
    \item Positive thresholds $\tau > 0$ induce stable decision regions with nonzero angular margin
    \item The zero threshold $\tau = 0$ yields a degenerate boundary
\end{itemize}

The canonical threshold $\tau = 1/\sqrt{2}$ (angle $\pi/4$) provides a symmetric reference point for balanced gating.

\end{document}